% ===============================================================
%  مقدمه: مرور پیش‌نیازها (جلسات ۱ و ۲)
% ===============================================================
\chapter*{پیش‌درآمد: مرور پیش‌نیازها (جلسات اول و دوم)}
\addcontentsline{toc}{chapter}{پیش‌درآمد: مرور پیش‌نیازها (جلسات اول و دوم)}
\markboth{پیش‌درآمد: مرور پیش‌نیازها}{}

\begin{note}
\notetitle این جزوه بر اساس اسلایدهای جلسه سوم به بعد نوشته شده است. از آنجا که در جلسات بعدی مدام به مفاهیم جلسات اول و دوم (نظریه مجموعه‌ها، اصول شمارش و تعریف فضای احتمال) ارجاع داده می‌شود، در این پیش‌درآمد خلاصه‌ای کامل و خودبسنده از این پیش‌نیازها آورده‌ایم تا جزوه به تنهایی قابل مطالعه باشد.
\end{note}

\section*{۱. نظریه مجموعه‌ها}
\addcontentsline{toc}{section}{نظریه مجموعه‌ها}

\textbf{مجموعه} گردایه‌ای از اشیای متمایز است که به آن‌ها \textbf{عضو} می‌گوییم. اگر $x$ عضو مجموعه $A$ باشد می‌نویسیم $x\in A$ و در غیر این صورت $x\notin A$. مجموعه‌ای که هیچ عضوی ندارد را \textbf{مجموعه تهی} می‌نامیم و با $\varnothing$ نشان می‌دهیم. اگر هر عضو $A$ عضو $B$ نیز باشد، می‌گوییم $A$ زیرمجموعه $B$ است و می‌نویسیم $A\subseteq B$.

در احتمال، همیشه یک «مجموعه مرجع» یا \textbf{فضای نمونه} $\Omega$ داریم و همه مجموعه‌ها زیرمجموعه آن هستند. عملگرهای اصلی روی مجموعه‌ها عبارت‌اند از:
\begin{itemize}
\item \textbf{متمم:} $A^c=\{\omega\in\Omega:\ \omega\notin A\}$.
\item \textbf{اجتماع:} $A\cup B=\{\omega:\ \omega\in A \ \text{or}\ \omega\in B\}$ (رخ دادن «$A$ یا $B$»).
\item \textbf{اشتراک:} $A\cap B=\{\omega:\ \omega\in A \ \text{and}\ \omega\in B\}$ (رخ دادن «$A$ و $B$»). گاهی $A\cap B$ را به صورت $AB$ یا $A,B$ هم می‌نویسیم.
\item \textbf{تفاضل:} $A\setminus B=A\cap B^c$ (رخ دادن $A$ بدون $B$).
\item \textbf{ضرب دکارتی:} $A\times B=\{(a,b):\ a\in A,\ b\in B\}$؛ برای مثال فضای نمونه پرتاب دو تاس برابر است با $\{1,\dots,6\}\times\{1,\dots,6\}$.
\end{itemize}

دو مجموعه $A$ و $B$ را \textbf{ناسازگار} (\lr{disjoint}) گوییم اگر $A\cap B=\varnothing$؛ یعنی هیچ‌گاه هم‌زمان رخ ندهند.

\begin{theorem}[قوانین دمورگان]
برای هر خانواده از مجموعه‌ها $\{A_i\}$:
\[
\Big(\bigcup_i A_i\Big)^c=\bigcap_i A_i^c,\qquad \Big(\bigcap_i A_i\Big)^c=\bigcup_i A_i^c .
\]
\end{theorem}
\begin{proof}
$\omega\in(\bigcup_i A_i)^c$ یعنی $\omega$ در هیچ‌یک از $A_i$ها نیست؛ یعنی برای هر $i$ داریم $\omega\in A_i^c$؛ یعنی $\omega\in\bigcap_i A_i^c$. رابطه دوم با اعمال رابطه اول روی $A_i^c$ها و گرفتن متمم از دو طرف به دست می‌آید.
\end{proof}

\section*{۲. اصول شمارش}
\addcontentsline{toc}{section}{اصول شمارش}

\begin{theorem}[اصل ضرب]
اگر انجام یک کار شامل $k$ مرحله باشد و مرحله $i$ام به $n_i$ طریق (مستقل از انتخاب‌های مراحل قبل) انجام شود، کل کار به $n_1n_2\cdots n_k$ طریق انجام می‌شود. معادلاً $|A_1\times\cdots\times A_k|=|A_1|\cdots|A_k|$.
\end{theorem}

\begin{itemize}
\item \textbf{جایگشت:} تعداد ترتیب‌های $n$ شیء متمایز برابر $n!$ است. تعداد ترتیب‌های $k$تایی از $n$ شیء برابر $\frac{n!}{(n-k)!}$ است.
\item \textbf{ترکیب:} تعداد زیرمجموعه‌های $k$عضوی یک مجموعه $n$عضوی برابر است با
\[\binom{n}{k}=\frac{n!}{k!\,(n-k)!}.\]
\item \textbf{قضیه دوجمله‌ای:} $(x+y)^n=\sum_{k=0}^{n}\binom{n}{k}x^ky^{n-k}$؛ به‌ویژه $\sum_{k=0}^n\binom nk=2^n$ که همان تعداد کل زیرمجموعه‌های یک مجموعه $n$عضوی است.
\end{itemize}

\begin{theorem}[ستاره و میله \lr{(Stars and Bars)}]\label{thm:stars}
تعداد جواب‌های معادله $x_1+x_2+\cdots+x_k=n$ در اعداد صحیح نامنفی برابر است با
\[\binom{n+k-1}{k-1}.\]
\end{theorem}
\begin{proof}
هر جواب را می‌توان با یک رشته از $n$ ستاره و $k-1$ میله نشان داد: تعداد ستاره‌های قبل از اولین میله برابر $x_1$، بین میله اول و دوم برابر $x_2$ و \ldots\ است. برای مثال با $k=3$ و $n=5$، رشته $\star\star|\star|\star\star$ نمایانگر جواب $(2,1,2)$ است. این تناظر یک‌به‌یک است، پس کافی است تعداد راه‌های انتخاب جای $k-1$ میله از میان $n+k-1$ جایگاه را بشماریم که برابر $\binom{n+k-1}{k-1}$ است.
\end{proof}

\begin{theorem}[اصل شمول و عدم شمول برای اندازه مجموعه‌ها]
برای مجموعه‌های متناهی:
\[|A\cup B|=|A|+|B|-|A\cap B|,\]
\[|A\cup B\cup C|=|A|+|B|+|C|-|A\cap B|-|A\cap C|-|B\cap C|+|A\cap B\cap C|.\]
\end{theorem}

\section*{۳. فضای احتمال و اصول کولموگروف}
\addcontentsline{toc}{section}{فضای احتمال و اصول کولموگروف}

\begin{definition}[فضای احتمال]
یک \textbf{فضای احتمال} یک سه‌تایی مرتب $(\Omega,\F,\Prob)$ است که در آن:
\begin{itemize}
\item $\Omega$ \textbf{فضای نمونه}، یعنی مجموعه همه نتایج ممکن یک آزمایش تصادفی است.
\item $\F$ یک \textbf{میدان سیگما} (\lr{$\sigma$-field}) روی $\Omega$ است؛ یعنی گردایه‌ای از زیرمجموعه‌های $\Omega$ (که به آن‌ها \textbf{رویداد} یا \textbf{پیشامد} می‌گوییم) با این ویژگی‌ها:
 (الف) $\Omega\in\F$؛ (ب) اگر $A\in\F$ آنگاه $A^c\in\F$؛ (ج) اگر $A_1,A_2,\ldots\in\F$ آنگاه $\bigcup_{i=1}^\infty A_i\in\F$.
\item $\Prob:\F\to[0,1]$ یک \textbf{اندازه احتمال} است که در اصول کولموگروف صدق می‌کند.
\end{itemize}
\end{definition}

\begin{definition}[اصول کولموگروف]
تابع $\Prob$ روی $\F$ یک اندازه احتمال است اگر:
\begin{enumerate}
\item \textbf{نامنفی بودن:} برای هر $A\in\F$، $\Prob(A)\ge 0$.
\item \textbf{نرمال بودن:} $\Prob(\Omega)=1$.
\item \textbf{جمع‌پذیری شمارا:} برای هر دنباله از رویدادهای دوبه‌دو ناسازگار $A_1,A_2,\ldots$:
\[\Prob\Big(\bigcup_{i=1}^{\infty}A_i\Big)=\sum_{i=1}^{\infty}\Prob(A_i).\]
\end{enumerate}
\end{definition}

به زبان ساده، اصول کولموگروف می‌گویند $\Prob$ یک «اندازه» (\lr{measure}) روی $\Omega$ است ــ مانند طول، مساحت یا تعداد اعضا ــ با این قید اضافه که اندازه کل فضا برابر ۱ باشد.

\begin{note}
\notetitle \textbf{تعبیر فرکانسی احتمال.} اگر آزمایشی را $N$ بار تکرار کنیم و رویداد $A$ در $N_A$ بار رخ دهد، فرکانس نسبی $N_A/N$ برای $N$های بزرگ به $\Prob(A)$ نزدیک می‌شود. فرکانس نسبی به‌وضوح نامنفی است، برای $\Omega$ برابر ۱ است و برای دو رویداد ناسازگار داریم $N_{A\cup B}=N_A+N_B$؛ پس اصول کولموگروف دقیقاً همان ویژگی‌های شهودی فرکانس نسبی هستند.
\end{note}

\begin{note}
\notetitle وقتی $\Omega$ متناهی یا شماراست، معمولاً $\F=2^\Omega$ (مجموعه همه زیرمجموعه‌های $\Omega$) انتخاب می‌شود. اگر $|\Omega|=n$ آنگاه $|2^\Omega|=2^n$.
\end{note}
